\section{Benchmarks}
\label{sec:benchmarks}

A software paper for an inference framework needs evidence that the framework is fast enough to be useful and that its scaling behaviour is well-characterised. Rather than report a handful of ad-hoc timings, we define a standardised ten-experiment benchmark suite for on-demand volumetric inference, spanning interactive latency, blockwise throughput, cluster scaling (strong and weak), memory efficiency, comparison against a hand-rolled baseline, and the cost of optional features (multi-model fusion, half-precision, alternative I/O backends, model-format adapters, and long-running stability). The full suite is summarised in Table~\ref{tab:benchmarks}. We execute three of the ten here --- B1, B3, and B6, the smallest set that establishes the core claims: interactive use is fast enough to be useful, whole-volume use scales out on a cluster, and the end-to-end story improves on the hand-rolled workflow it replaces. The remaining seven are specified as fully detailed, reproducible protocols so that they can be run on other hardware and datasets and contributed back. Each benchmark below is described in enough detail to be reproduced independently.

\begin{table*}[tbhp]
\centering
\caption{Standardised benchmark suite for on-demand volumetric inference. ``Status'' marks the three experiments executed in this manuscript, each with a pointer to its results table; the remaining seven are specified as reproducible protocols in the text.}
\label{tab:benchmarks}
\small
\begin{tabularx}{\textwidth}{l X l l}
\toprule
\textbf{ID} & \textbf{Question} & \textbf{Primary metric} & \textbf{Status} \\
\midrule
B1 & Interactive chunk latency vs.\ chunk size and model & ms / chunk (median, p99) & executed (Tab.~\ref{tab:b1_results}) \\
B2 & Blockwise throughput on a single GPU & GVox/s & protocol \\
B3 & Strong scaling on cluster (fixed volume, varying GPUs) & speedup vs.\ ideal & executed (Tab.~\ref{tab:b3_results}) \\
B4 & Weak scaling on cluster (fixed work per GPU) & efficiency & protocol \\
B5 & Memory footprint vs.\ chunk size (with/without FP16) & peak GPU/CPU MB & protocol \\
B6 & Comparison vs.\ hand-rolled PyTorch + offline export & end-to-end wall time & executed (Tab.~\ref{tab:b6_results}) \\
B7 & Multi-model fusion overhead (1 vs.\ 2 vs.\ 3 models) & relative latency & protocol \\
B8 & I/O backend comparison (Zarr / N5 / OME-Zarr / S3 / precomputed) & read latency & protocol \\
B9 & Model-format adapter cost (TorchScript / DaCapo / BioImage.io / HuggingFace) & load + first-chunk time & protocol \\
B10 & Long-running stability (24~h server, real workload) & failure rate, drift & protocol \\
\bottomrule
\end{tabularx}
\end{table*}

\subsection*{B1: Interactive chunk latency}

\paragraph{Question.} How quickly does \cellmapflow{} return a chunk in server mode, as a function of chunk size and model size? This is the single most user-visible metric: it determines whether scrolling through a Neuroglancer view feels instantaneous or laggy.

\paragraph{Protocol.} Use two published OpenOrganelle datasets that span the cell-culture / tissue axis and a meaningful size disparity: \texttt{jrc\_hela-2} (HeLa cell line, $\sim$122~GB at \texttt{s0}) and \texttt{jrc\_mus-liver} (mouse liver, $\sim$1.45~TB at \texttt{s0}, $\sim$12$\times$ larger), both from Heinrich et al.\ 2021 \cite{Heinrich2021}. The size disparity is deliberate: \texttt{jrc\_hela-2} still fits in main memory, while \texttt{jrc\_mus-liver} does not, so the larger dataset exercises the chunked-inference value proposition that motivates \cellmapflow{}. The medium-size model is \texttt{cellmap/fly\_organelles\_run07\_700000}, a 3D U-Net with native input shape $178^3$ at 8~nm voxel size. For each dataset, request 200 chunks at random spatial offsets through the HTTP server after a 20-chunk warm-up; report median, p95, and p99 wall-clock time per chunk.

\paragraph{First-pass results.} On a single NVIDIA H100 80GB GPU reading inputs from Janelia's \texttt{/nrs/cellmap} mount, both datasets land squarely in interactive territory: chunk-request latencies are well under 250~ms at p99 in both cases (Tab.~\ref{tab:b1_results}, Fig.~\ref{fig:b1_latency}). The tissue dataset is $\sim$70\% slower per chunk than the cell-culture dataset, consistent with denser content compressing less and producing larger Zarr chunks to read; chunk payload sizes range from $\sim$700~bytes to $\sim$5~MB across the random-offset sweep on \texttt{jrc\_hela-2}.

\begin{figure}[tbhp]
\centering
\includegraphics[width=\linewidth]{b1_latency}
\caption{B1 interactive chunk-request latency on a single H100, over 200 random-offset requests per dataset after a 20-request warm-up. Boxes span the interquartile range, the red line is the median, whiskers the 5th/95th percentiles, and the diamond marks p99. Both datasets stay under 250~ms at p99; the tissue volume (\texttt{jrc\_mus-liver}) is slower and more variable than the cell-culture volume (\texttt{jrc\_hela-2}).}
\label{fig:b1_latency}
\end{figure}

\paragraph{Future sweeps.} The full sweep promised here --- chunk sizes $\{64^3, 128^3, 256^3\}$ $\times$ three model sizes (small UNet $\sim$1M params, medium $\sim$10M, large $\sim$50M) $\times$ \{H100, A100, V100, CPU\} --- remains to be executed. The first-pass results above establish the medium-cell of that grid.

\subsection*{B2: Blockwise throughput on a single GPU}

\paragraph{Question.} How many output voxels per second does \cellmapflow{} produce in blockwise mode on a single GPU?

\paragraph{Protocol.} Run \texttt{cellmap\_flow\_blockwise} on a fixed-size sub-volume (e.g.\ $4096^3$~voxels) with one worker on H100. Vary chunk size and FP16 on/off. Report sustained GVox/s, GPU utilisation, and time spent in I/O versus compute (instrument the inferencer's load/forward/postprocess phases).

\paragraph{Expected outcome.} A baseline against which multi-GPU runs are compared in B3/B4. Useful side product: a recommendation for default chunk size as a function of model and hardware.

\subsection*{B3: Strong scaling on cluster}

\paragraph{Question.} For a fixed-size volume, how does total wall time scale with the number of GPU workers?

\paragraph{Protocol.} Run on both \texttt{jrc\_hela-2} (cell culture) and \texttt{jrc\_mus-liver} (tissue) with the same model. Run \texttt{cellmap\_flow\_blockwise} with worker counts $N \in \{1, 4, 16, 64, 128\}$ on the LSF cluster. Hold all other parameters fixed within each dataset. Report wall time, speedup relative to $N=1$, parallel efficiency, and number of Daisy retries. Plot the scaling curve per dataset and fit Amdahl's-law parameters; comparing the two curves separates dataset-specific effects (chunk-content variance, S3 read patterns) from the framework's intrinsic scaling behaviour.

\paragraph{Results.} We ran the full sweep on both datasets over a fixed $8192^3$~nm sub-volume ($\sim$6{,}900 output blocks at the model's $56^3$ write size), with \texttt{cellmap\_flow\_blockwise} dispatching $N$ single-GPU H100 workers via Daisy (Tab.~\ref{tab:b3_results}, Fig.~\ref{fig:b3_scaling}). On \texttt{jrc\_hela-2}, wall time falls from 1292~s at $N{=}1$ to 122~s at $N{=}128$ (a $10.5\times$ speedup); on \texttt{jrc\_mus-liver}, from 1669~s to 132~s ($12.7\times$). Scaling is near-linear through $N{=}4$ (parallel efficiency 0.80--0.86) and then bends over as fixed per-worker costs --- LSF dispatch, GPU allocation, and model load --- grow relative to each worker's shrinking share of compute; efficiency falls to $\sim$0.1 by $N{=}128$, although wall time still improves at every step. The tissue volume scales consistently better than the cell-culture volume (0.36 vs.\ 0.30 efficiency at $N{=}16$; $12.7\times$ vs.\ $10.5\times$ at $N{=}128$): its heavier per-block compute and I/O amortise the fixed per-worker overhead more effectively, so the framework extracts more parallelism from the larger workload. The practical takeaway is that a few to a few dozen workers captures most of the available speedup for a sub-volume of this size, with diminishing but still-positive returns beyond --- an empirical estimate of \cellmapflow{}'s serial overhead and the useful worker count per dataset.

\begin{figure}[tbhp]
\centering
\includegraphics[width=\linewidth]{b3_scaling}
\caption{B3 cluster strong scaling on a fixed $8192^3$~nm sub-volume, with \texttt{cellmap\_flow\_blockwise} dispatching $N$ single-GPU H100 workers via Daisy. Left: wall time vs.\ worker count (log--log; dashed line is ideal linear scaling from $N{=}1$). Right: parallel efficiency. Scaling tracks the ideal through $N{=}4$, then departs as per-worker fixed costs (scheduler dispatch, GPU allocation, model load) come to dominate each worker's shrinking compute share.}
\label{fig:b3_scaling}
\end{figure}

\subsection*{B4: Weak scaling on cluster}

\paragraph{Question.} If the work per GPU is held constant, does adding GPUs maintain the same wall time?

\paragraph{Protocol.} Hold per-GPU work fixed at, e.g., $1024^3$ voxels per worker. Run with $N \in \{1, 4, 16, 64\}$ workers, each processing a disjoint sub-volume of the dataset. Report wall time per worker, parallel efficiency, and any tail-latency degradation.

\paragraph{Expected outcome.} Identifies whether scheduling, file-system, or network bottlenecks degrade performance as the worker pool grows, separately from inherent serial fraction.

\subsection*{B5: Memory footprint}

\paragraph{Question.} What are peak GPU and CPU memory consumption as a function of chunk size and FP16 use?

\paragraph{Protocol.} Profile a single worker over a representative pipeline. Sweep chunk size $\{64^3, 128^3, 192^3, 256^3\}$ and FP16 on/off. Capture peak \texttt{torch.cuda.max\_memory\_allocated}, peak resident-set size (RSS), and any out-of-memory failures.

\paragraph{Expected outcome.} A practical guide for users picking chunk size on a given GPU. FP16 should approximately halve GPU memory and modestly reduce latency; the experiment quantifies the trade-off.

\subsection*{B6: Comparison against a hand-rolled baseline}

\paragraph{Question.} What does \cellmapflow{} actually save a user compared with the workflow it replaces?

\paragraph{Protocol.} Define a representative task: ``produce viewable predictions for organelle X across dataset Y.'' Implement two pipelines:
\begin{itemize}
\item \emph{Baseline.} A hand-written PyTorch script that loops over chunks, runs inference, writes output to Zarr, then registers the result with Neuroglancer's precomputed adapter.
\item \emph{\cellmapflow{}.} The same model + dataset configured via YAML, server mode for the first qualitative look, blockwise mode for full-volume export.
\end{itemize}
Measure two end-to-end metrics: (a) \emph{time to first interactive view} (from cold start to the user seeing the first prediction); (b) \emph{time to whole-volume completion}. Track lines of code as a secondary metric.

\paragraph{Results.} We ran both pipelines on identical hardware (one H100 per worker) over the same $8192^3$~nm sub-volume used in B3, on both a cell-culture (\texttt{jrc\_hela-2}) and a tissue (\texttt{jrc\_mus-liver}) volume, with \texttt{cellmap/fly\_organelles\_run07\_700000} (Tab.~\ref{tab:b6_results}). Time-to-first-view is interactive either way: a single chunk returns in seconds (7--17~s, dominated by model load). The completion story turns on parallelism. At a single worker, \cellmapflow{} is about $2\times$ slower than the hand-rolled baseline (1324 vs.\ 608~s on \texttt{jrc\_hela-2}; 1691 vs.\ 790~s on \texttt{jrc\_mus-liver}), because its blockwise path routes every block through Daisy and the LSF scheduler, which at $N{=}1$ is pure overhead. But the same YAML scales out with no code change: at $N{=}4$ workers \cellmapflow{} already overtakes the single-process baseline on both datasets, and at $N{=}16$ it finishes roughly $5$--$6\times$ faster (122 vs.\ 608~s; 135 vs.\ 790~s) --- an $11$--$12\times$ speedup over its own single-worker time, tracking the near-linear-then-saturating scaling characterised in B3. The single-process baseline cannot follow: parallelising it is exactly the chunked, scheduled, restartable machinery the user would otherwise have to build by hand.

\paragraph{What \cellmapflow{} actually saves.} The operational difference compounds the throughput one. The baseline is a 114-line hand-written inference script (chunked reads, context padding, model invocation, Zarr writes, edge handling) that the user must write, debug, and maintain; the \cellmapflow{} run is a 16-line YAML config plus one CLI invocation, with the same chunked inference, multi-format I/O, viewer integration, and multi-worker scale-out all handled by the framework --- a $7\times$ reduction in user-written code. So the honest framing of B6 is: comparable interactive latency; a single-worker completion penalty that is the price of the scheduling layer; and, once that layer is used as intended, both a decisive completion-time win over the hand-rolled baseline and $7\times$ less code to get there.

\paragraph{Caveats.} These are single-run timings (no repeats or error bars) for one model on one $8192^3$~nm sub-volume of each dataset, on H100s. Parallel efficiency falls as $N$ grows (B3), so the $N{=}16$ advantage is a throughput result rather than a claim of ideal scaling, and the crossover point ($N{\approx}4$ here) will shift with sub-volume size, model cost, and I/O backend.

\subsection*{B7: Multi-model fusion overhead}

\paragraph{Question.} How much does adding a second or third model to an inference pipeline cost?

\paragraph{Protocol.} Run B1 with one, two, and three models loaded simultaneously, fused with each merge mode (\texttt{AND}, \texttt{OR}, consensus). Measure incremental latency per added model and peak GPU memory.

\paragraph{Expected outcome.} A user-facing rule of thumb: ``adding model $k$ costs approximately $X\%$ more latency.''

\subsection*{B8: I/O backend comparison}

\paragraph{Question.} How does input-data backend affect read latency?

\paragraph{Protocol.} Mirror the same dataset to local Zarr, local N5, local OME-Zarr, S3-hosted Zarr (cold and warmed CDN), and Google Cloud precomputed. Sweep B1's chunk-request workload over each backend. Report read latency only (skip inference, or subtract inference time using B1 numbers).

\paragraph{Expected outcome.} Quantifies the cloud-vs.-local penalty and validates the TensorStore abstraction. Will inform recommended deployment patterns.

\subsection*{B9: Model-format adapter cost}

\paragraph{Question.} Do the various supported model formats incur different load and first-inference times?

\paragraph{Protocol.} For one logical model, package it as TorchScript, DaCapo, BioImage.io, and HuggingFace Hub variants. Measure cold-start load time and first-chunk inference time per format. Subsequent-chunk latency should be format-independent; verify.

\paragraph{Expected outcome.} A guide to format choice when trade-offs matter (e.g.\ TorchScript for shortest cold start, BioImage.io for portability).

\subsection*{B10: Long-running stability}

\paragraph{Question.} Does the server hold up under realistic continuous use?

\paragraph{Protocol.} Run a server for 24~hours under a synthetic Neuroglancer-like workload (chunk requests at 1--10~Hz, varied spatial pattern). Track failure events, GPU memory growth, response-time drift, and cache effectiveness. Repeat in blockwise mode for a multi-hour batch run.

\paragraph{Expected outcome.} Quantitative evidence of production readiness, or a list of stability bugs to fix before publication.

\subsection*{Recommended order of execution}

The suite is designed to be run incrementally, and we suggest an order for extending beyond the three experiments executed here. The core \textbf{B1, B3, and B6} come first: B1 establishes that interactive use is fast enough to be useful; B3 establishes that whole-volume use scales out; B6 establishes that the framework's end-to-end story beats the baseline it replaces. B5 (memory) and B9 (model-format adapter cost) are quick to add once the harness from B1 exists. B7 and B8 are most valuable as supporting evidence for the discussion. B4 and B10 are the most involved and chiefly strengthen production-readiness claims.

\subsection*{Reporting and reproducibility}

The three executed benchmarks (B1, B3, B6) are implemented as runnable scripts in the \texttt{benchmarks/} directory accompanying this manuscript, with version pins for the framework, models, and datasets; per-experiment YAML configs are checked in alongside their results so that any third party can reproduce the corresponding tables. The remaining seven are specified by the protocols above in enough detail --- inputs, sweeps, primary metrics, and expected outcomes --- to be implemented and run independently on other hardware and datasets. Hardware (GPU model, interconnect, file-system class) is reported with each result. Benchmark inputs are drawn from the OpenOrganelle FIB-SEM release \cite{Xu2021,OpenOrganelle}; trained models are obtained from runs against the CellMap Segmentation Challenge dataset \cite{CellMapChallenge} via DaCapo \cite{Patton2024,DaCapoRepo}.
